\documentclass[11pt,a4paper]{article}
\usepackage[T1]{fontenc}
\usepackage{lmodern,amsmath,amssymb,amsthm,mathtools,booktabs,microtype}
\usepackage[margin=25mm]{geometry}
\usepackage[colorlinks=true,linkcolor=blue,urlcolor=blue,citecolor=blue]{hyperref}
\setlength{\emergencystretch}{3em}
\newtheorem{theorem}{Theorem}[section]
\newtheorem{lemma}[theorem]{Lemma}
\newtheorem{proposition}[theorem]{Proposition}
\newcommand{\Q}{\mathbb Q}\newcommand{\Z}{\mathbb Z}
\DeclareMathOperator{\GL}{GL}
\title{Primitive quadratic bases and small-coefficient quartic presentations}
\author{}\date{15 November 2026}
\begin{document}\maketitle
\begin{abstract}
Invertible rational changes of basis among quadratic equations give
small-coefficient presentations of two defining systems. The first has
196 total variables, 176 primitive quadratic summands and 2994 monomials;
the second has 78 total variables, 68 primitive quadratic summands and
1254 monomials. We give exact coefficient-height formulas for every allowed
coefficient parameter. The rational zero sets are preserved, up to two
specified nonzero constant coordinate scalings in the first family.
No arithmetic premise or sufficient defining parameter is changed.
The changes improve coefficients without reducing the variable counts.
\end{abstract}

\section{The two quartic families}
A variable count includes the distinguished rational input $X$.
A coefficient parameter is fixed before existential quantification.
For an integer polynomial $F$, write $H(F)$ for the largest absolute
value of its integer coefficients after the parameter is fixed.
Set
\begin{align*}
 K_0&=126616122944180358034867798936,\\
 K_2&=32561227748626372080,\\
 J_0&=20828275368498373068800,\qquad
 J_1=275933278576697172.
\end{align*}

\begin{theorem}\label{thm:families}
There are complete integer coefficient specifications for families
$P_N\in\Z[X,Y_1,\ldots,Y_{195}]$ and
$Q_B\in\Z[X,Z_1,\ldots,Z_{77}]$ with the following properties.
\begin{enumerate}
\item For every integer $N\ge1$, $P_N$ is a sum of 176 primitive
quadratic squares, has total degree exactly four, integer content one,
and exactly 2994 coordinate monomials. Its exact height is
\[
                    H(P_N)=\max\{K_0,K_2\,149^{2N}\}.
\]
The first term dominates for $N=1,2$, and the second for $N\ge3$.
Its rational input projection equals that of the 196-variable parent
of~\cite{Parent196} for every fixed $N$.
\item For every positive integer $B$, $Q_B$ is a sum of 68 primitive
quadratic squares, has total degree exactly four, integer content one,
and exactly 1254 coordinate monomials. Its exact height is
\[
                         H(Q_B)=\max\{J_0,B^2+J_1\}.
\]
Its whole rational zero set equals that of the 68-row paired-ratio
system of~\cite{Paired}, for every fixed $B$.
\end{enumerate}
Consequently, under the precise inherited arithmetic premises of the
respective parent constructions, $P_N$ defines $\Z$ over $\Q$ for every
sufficiently large fixed $N$, and $Q_B$ does so for every sufficiently
large fixed odd $B$. The latter assertion retains ordinary rational
$abc$ as an additional hypothesis. Both sufficient thresholds are the
parent thresholds and remain ineffective.
\end{theorem}

The $N=3$ comparison in the first assertion concerns coefficient size;
it is not a numerical defining threshold. In the second assertion the
algebraic statements hold for every positive $B$, while the inherited
arithmetic theorem uses sufficiently large odd $B\ge81$.
The 196-variable source retains the original attribution of the supplied
Conant--ChatGPT 198-variable construction~\cite{Archive}.
The full external arithmetic audit is not supplied by the finite
coefficient certificates proved here.

\section{Changing a basis of quadratic equations}
Let $R=\Z[\tau,x_1,\ldots,x_n]$, with the coefficient parameter $\tau$
assigned degree zero when taking degrees in the $x_i$. A row is primitive
if the greatest common divisor of its integer coefficients is one.

\begin{lemma}\label{lem:basis}
Let $q_1,\ldots,q_r\in\Q[\tau,\boldsymbol x]$ have degree at most two
in $\boldsymbol x$. Let $M\in\GL_r(\Q)$, and let
$x_i=d_i y_i$ with $d_i\in\Q^\times$, fixing the input coordinate.
Assume every row of $M\boldsymbol q(\tau,d_1y_1,\ldots,d_ny_n)$ is
nonzero, and multiply each by a nonzero rational constant to obtain a
primitive integer row. The resulting rows have degree at
most two and define a rational zero set in input-preserving bijection
with the original one, at every fixed rational value of $\tau$.
Their sum of squares has the same rational input projection.
\end{lemma}
\begin{proof}
The coordinate change is invertible over $\Q$, independently of
$\tau$. The vector of transformed rows vanishes if and only if the
original vector does, because $M$ has a rational inverse. Multiplying
individual equations by nonzero rational constants does not alter their
zero set. Each coordinate substitution is linear, so the degree bound
is preserved. Over $\Q$ a sum of squares vanishes exactly when all its
summands vanish. These statements also allow an unchanged collection
of further equations to be adjoined.
\end{proof}

There is no requirement that $M$ be integral or unimodular over $\Z$.
If $W$ is the rational span of coefficient vectors of some independent
quadratic rows, its primitive integer vectors lie in
$W\cap\Z^m$, where $m$ is the number of coefficient positions.
A basis of this lattice may be strictly smaller in coefficient size
than every integral basis change of the initially given row lattice.
For example, an integer decomposition $UAV=D$ of a full-row-rank
coefficient matrix into Smith form gives a rational basis change
\[
 R=\operatorname{diag}(D_{11}^{-1},\ldots,D_{rr}^{-1})U,
 \qquad RA\in\operatorname{Mat}_{r\times m}(\Z).
\]
Indeed $(RA)V=(I_r\;0)$, and $V^{-1}$ is integral. Subsequent invertible
integer row changes and primitive normalization remain covered by
Lemma~\ref{lem:basis}. The specifications below use fixed rational
matrices and their exact inverses. Their correctness requires no claim
of an optimal basis and no search during reconstruction.

\section{The 196-variable coefficient specification}
Use the balanced 176-row system of~\cite{Parent196}, with $C=149^N$.
Its source digest and all coordinate names are fixed by
\begin{center}\small
\nolinkurl{computation/quartic/primitive_basis196_v40/transform.json}.
\end{center}
Apply its displayed rational $8\times8$ matrix to the two rows
\texttt{testSquare.40.0.0.definition},
\texttt{testSquare.40.0.1.definition}, followed in order by the six
\texttt{sharedT.reference.relation}$i$.$j$ rows, with
$i=1,2,3$ and $j=0,1$. Its determinant is
\[
             \frac1{214450708242592116705963714180667354532074632082}\ne0.
\]
All other rows are initially unchanged. Next, for each of the two
coordinates named \texttt{additive.detector}$i$\texttt{.contactCombination},
$i=0,1$, make the constant change
\[
                   x_{\rm parent}=\frac{x_{\rm new}}{3813373161796}.
\]
These are coordinate indices 76 and 90. All other coordinates, including
$X$, are fixed. Clear every row to primitive integer coefficients.
The fixed matrix, coordinate map and primitive row multipliers specify
all 176 new quadratic equations without an optimization algorithm.
Primitivity also holds after every specialization $C=149^N$.
The checker verifies that either the parameter-independent coefficients
of a row already have greatest common divisor one, or each of its
coefficient profiles is a single power of $C$ and the greatest common
divisor of its parameter-independent coefficients is coprime to 149.
In the second case no prime dividing $C$ can divide every specialized
coefficient, and any other common prime would divide every formal
coefficient, contrary to the checked formal primitivity.

The independent checker constructs the rational inverse matrix and
verifies every transformed coefficient. The resulting map is a bijection
of rational solutions by Lemma~\ref{lem:basis}; the inverse first undoes
the two constant coordinate scalings and then uses the equality of row
spans. Thus no completeness witness or exceptional point is discarded.

The complete quadratic and quartic tables are \texttt{quadratic.json}
and \texttt{polynomial.json} in the same directory. For a coordinate
monomial $m$, let its coefficient in the quartic be
\[
                         p_m(C)=a_{m,0}+a_{m,1}C+a_{m,2}C^2.
\]
The exact finite certificate checks these assertions:
\begin{enumerate}
\item There are 2994 such monomials, and if $d=\max\{j:a_{m,j}\ne0\}$,
\[
             |a_{m,d}|149^d>\sum_{j<d}|a_{m,j}|149^j.
\]
Hence none disappears for $C\ge149$.
\item For every $m$,
\[
 \frac{|a_{m,0}|}{149^6}+\frac{|a_{m,1}|}{149^3}+|a_{m,2}|\le K_2.
\]
Thus $|p_m(C)|\le K_2C^2$ for all $C\ge149^3$.
\item The monomial of coordinate indices $(15,18,152,152)$ has coefficient
$-K_2C^2$. The leading bound is therefore attained.
\item Direct exact substitution at $C=149$ and $C=149^2$ gives height
$K_0$, and $K_2 149^4<K_0<K_2 149^6$.
\end{enumerate}
These finite integer inequalities prove the exact height formula for
every $N\ge1$, rather than just for sampled values. The input fourth
power still has coefficient one, as its defining square row is unchanged
and no other row contains $X$. This proves the degree and content claims.

\section{The 78-variable coefficient specification}
Use the 68-row source from~\cite{Paired}, with fixed parameter $B$.
The file
\begin{center}\small
\nolinkurl{computation/quartic/abc78_balanced_v40/basis-certificate.json}
\end{center}
specifies an invertible rational $8\times8$ transformation of the
reference rows, in the order
\[
 \texttt{common.fSquare.0},\quad\texttt{common.fSquare.1},
 \quad\texttt{common.first.0},\quad\texttt{common.first.1},
\]
\[
 \texttt{common.second.0},\quad\texttt{common.second.1},
 \quad\texttt{common.third.0},\quad\texttt{common.third.1}.
\]
Its exact determinant is
\[
               \frac1{107225354121296058352981857090333677266037316041}\ne0.
\]
The certificate includes its rational inverse. All other rows and all
coordinates are unchanged. Primitive normalization and
Lemma~\ref{lem:basis} prove equality of the entire rational zero sets
for every fixed $B$.
For each row the parameter-independent coefficients have greatest
common divisor one, so every specialization at an integer $B$ is
primitive as well.

The complete tables \texttt{quadratic-family.json} and
\texttt{polynomial-family.json} have 68 rows and 1254 monomials,
respectively. The independent checker expands the rows separately and
checks every coefficient as a polynomial in $B$. The only coefficient
of degree two in $B$ is the constant term $B^2+J_1$. It proves that all
other coefficient magnitudes are at most
$\max\{J_0,B^2+J_1\}$ for every positive $B$, and identifies a fixed
coefficient with magnitude $J_0$. The two candidates are therefore
both attained. Its sign checks also establish that no supported
monomial disappears for positive $B$. The untouched input-square
row again supplies coefficient one at $X^4$.

This establishes every algebraic assertion of
Theorem~\ref{thm:families}. The defining assertions follow from the
corresponding parent input-projection theorems, with exactly their
arithmetic premises and sufficient parameters. In particular the
$B^2$ term remains part of the coefficient height; improving the fixed
coefficient part does not reduce that parameter contribution.

\section{Reconstruction and verification}
The producers apply the fixed matrices and specified coordinate scalings.
The independent checkers import no producer code. They authenticate the
parent coefficient files, check the rational inverses and all row
identities, and expand the quadratic squares by separate algorithms.
The 196-variable checker uses only standard-library exact integer and
fraction arithmetic; the 78-variable checker also uses exact symbolic
expansion. The certificates concern the algebraic equivalences and
coefficient formulas, and do not replace the external arithmetic proofs.

\begin{thebibliography}{9}
\bibitem{Parent196}\emph{A 196-variable quartic with 176 quadratic
summands}, project manuscript, 15 November 2026,
\nolinkurl{papers/quartic/quartic196_redundancy_20261115_v40.tex}.
\bibitem{Paired}\emph{Recovery of elliptic curve equations from complete
ratio equations}, project manuscript, 15 November 2026,
\nolinkurl{papers/elliptic_charts/paired_ratio_redundancy_20261115_v40.tex}.
\bibitem{ABC78}\emph{A quartic in 78 variables from a mixed-power height estimate},
project manuscript, 15 November 2026,
\nolinkurl{papers/quartic_existence/quartic_abc78_20261115_v34.tex}.
\bibitem{Laurent}\emph{Quartics in 196 and 258 variables by elimination
after monomial rescaling}, project manuscript, 15 November 2026,
\nolinkurl{papers/quartic/laurent_elimination_20261115_v38.tex}.
\bibitem{Archive}\emph{A quartic in 198 variables}, supplied research
archive, ChatGPT in collaboration with Gabriel Conant's research request,
14 November 2026,
\nolinkurl{sources/h10_quartic198/arithmetic/inverse_vieta_198.tex}.
\end{thebibliography}
\end{document}
